\section{Statistical Methodology}
\label{sec:appendix_inference}

\subsection{Fractional Factorial Designs and Aliasing}

Experiment~1 uses a fractional factorial design to reduce experimental burden. The key trade-off is aliasing, whereby some higher-order interactions become confounded with lower-order terms. Experiments~2 and~3 use mixed-level designs (combining the three-level $\eta$ factor with binary factors), and Experiments~4a and~4b each use a full $2^6$ factorial, so aliasing is not a concern for these experiments.

\subsubsection{Resolution~V (Experiment~1)}

In a Resolution~V\footnote{In a Resolution~$R$ design, no $p$-factor interaction is aliased with any interaction of fewer than $R - p$ factors. Resolution~V therefore guarantees that main effects are free of two-way interaction bias and two-way interactions are free of other two-way interaction bias.} design, no main effect is aliased with any interaction of fewer than four factors, and no two-way interaction is aliased with any interaction of fewer than three factors. This ensures that all main effects and two-way interactions are estimable without bias from three-way terms, assumed negligible under effect sparsity. We validate this assumption using LASSO variable selection and nonparametric model comparison.

\subsection{Inference Corrections}

We compute HC3 robust standard errors \citep{MacKinnonWhite1985} to account for non-constant error variance across the design space. Comparing OLS standard errors to HC3 robust standard errors, we verify that significance claims hold under heteroscedasticity. The fraction of effects changing significance status under HC3 ranges from 0\% (Experiment~4a) to 6.7\% (Experiment~1); per-experiment details appear in Section~\ref{sec:appendix_robustness}.

The hat matrix $H = X(X^\top X)^{-1}X^\top$ has diagonal elements $h_{ii}$ (leverages). The HC3 sandwich estimator is
\begin{align}
  \hat{V}_{\text{HC3}} = (X^\top X)^{-1} X^\top \operatorname{diag}\!\Bigl(\frac{\hat{e}_i^2}{(1-h_{ii})^2}\Bigr) X\, (X^\top X)^{-1},
  \label{eq:hc3}
\end{align}
where $\hat{e}_i$ denotes the $i$-th OLS residual.

With $k$ main effects, $\binom{k}{2}$ two-way interactions, and $m$ response variables, each experiment involves many simultaneous hypothesis tests. We apply Holm--Bonferroni sequential correction \citep{Holm1979} to control the family-wise error rate at $\alpha = 0.05$. Key findings, including the auction type main effect and the auction type $\times$ exploration interaction, survive this stringent correction across all responses.

The family-wise error rate is the probability of at least one false rejection, defined as $\text{FWER} = \Pr(V \geq 1) \leq \alpha$, where $V$ is the number of false rejections. Given ordered $p$-values $p_{(1)} \leq \cdots \leq p_{(m)}$, the Holm step-down procedure rejects $H_{(k)}$ while $p_{(k)} \leq \alpha/(m - k + 1)$, stopping at the first non-rejection. The adjusted $p$-values are $\tilde{p}_{(k)} = \max_{j \leq k}\{(m - j + 1)\,p_{(j)}\}$, ensuring monotonicity.

We also apply Benjamini--Hochberg (BH) correction \citep{BenjaminiHochberg1995} to control the false discovery rate (FDR), a less conservative criterion that is better suited to exploratory screening of many effects. The FDR is the expected proportion of false discoveries among rejections: $\text{FDR} = \mathbb{E}[V/R]$, where $V$ is the number of false rejections and $R = \max(R_{\text{obs}}, 1)$ is the total number of rejections. The BH step-up procedure finds the largest $k$ such that $p_{(k)} \leq (k/m)\,q$ and rejects $H_{(1)}, \ldots, H_{(k)}$. The FWER controls the probability of \emph{any} false positive; the FDR controls the \emph{rate} of false positives among discoveries, making it more powerful when many effects are truly active.

For the top~10 effects by absolute $t$-statistic in each response, we compute Rademacher wild bootstrap $p$-values (1{,}000 iterations) to validate inference under minimal distributional assumptions. Bootstrap $p$-values align closely with HC3 asymptotic $p$-values (differences below 0.01), confirming robustness.

The Rademacher wild bootstrap assigns random weights $w_i^* \in \{-1, +1\}$ with equal probability and constructs pseudo-data under the null hypothesis $H_0\!: \beta_j = 0$:
\begin{align}
  y_i^* = \mathbf{x}_i^\top \tilde{\boldsymbol{\beta}} + \tilde{e}_i \cdot w_i^*,
  \label{eq:wildboot}
\end{align}
where $\tilde{\boldsymbol{\beta}}$ and $\tilde{e}_i$ are obtained from the restricted model imposing the null. The bootstrap $p$-value is $\hat{p}^* = B^{-1}\sum_{b=1}^{B} \mathbf{1}(|t^{*b}| \geq |t|)$, where $t^{*b}$ is the $t$-statistic from the $b$-th bootstrap sample. Using restricted rather than unrestricted residuals ensures that the bootstrap distribution is centred under $H_0$, improving size accuracy in finite samples \citep{DavidsonFlachaire2008}.

We fit quantile regressions at the 10th, 25th, 50th, 75th, and 90th percentiles of each response to assess whether factor effects are uniform across the outcome distribution or concentrated in the tails. An effect that is large at extreme quantiles but small at the median indicates that the factor primarily influences worst-case or best-case configurations rather than the typical outcome. This complements the OLS analysis, which estimates effects at the conditional mean.

The quantile regression estimator \citep{KoenkerBassett1978} minimises the asymmetric check function:
\begin{align}
  \hat{\boldsymbol{\beta}}(\tau) = \arg\min_{\boldsymbol{\beta}} \sum_{i=1}^{n} \rho_\tau\!\bigl(y_i - \mathbf{x}_i^\top \boldsymbol{\beta}\bigr), \qquad \rho_\tau(u) = u\bigl(\tau - \mathbf{1}\{u < 0\}\bigr).
  \label{eq:quantreg}
\end{align}
The coefficient $\beta_j(\tau)$ measures the marginal effect of factor $j$ on the $\tau$-th conditional quantile of the response, allowing us to distinguish factors that shift the entire distribution from those that primarily affect the tails.

\subsubsection{Lenth's Pseudo-Standard Error}

For each response variable, we compute Lenth's pseudo-standard error \citep{Lenth1989}: $\text{PSE} = 1.5 \cdot \text{median}_{|e_j| < 2.5\,s_0} |e_j|$, where $s_0 = 1.5 \cdot \text{median}_j |e_j|$ and $e_j$ are the estimated effects (OLS coefficients excluding the intercept). The margin of error $\text{ME} = t_{\alpha/2,\, d} \cdot \text{PSE}$ (with $d = \lfloor k/3 \rfloor$) provides an individual significance threshold, and the simultaneous margin of error $\text{SME} = t_{\gamma,\, d} \cdot \text{PSE}$ (where $\gamma = 1 - \alpha/(2k)$) applies a Bonferroni correction. Effects exceeding ME on the half-normal plot are flagged as likely active. In our replicated designs, Lenth's method serves as a complementary check rather than the primary inference tool, since we have independent error estimates from replication.

\subsection{Model Adequacy}

With two replicates per cell, we decompose the residual sum of squares into pure error (within-cell variation) and lack of fit (deviation of cell means from model predictions). The $F$-test $\text{MS}_{\text{LOF}} / \text{MS}_{\text{PE}}$ assesses whether the linear model with two-way interactions adequately fits the data. Per-experiment lack-of-fit results are reported in Section~\ref{sec:appendix_robustness}.

Formally, the residual sum of squares decomposes as $\text{SS}_{\text{Res}} = \text{SS}_{\text{LOF}} + \text{SS}_{\text{PE}}$, where
\begin{align}
  \text{SS}_{\text{PE}} &= \sum_{i=1}^{c} \sum_{j=1}^{n_i} (y_{ij} - \bar{y}_{i\cdot})^2, &\quad \text{df} &= n - c, \notag \\
  \text{SS}_{\text{LOF}} &= \sum_{i=1}^{c} n_i (\bar{y}_{i\cdot} - \hat{y}_i)^2, &\quad \text{df} &= c - p,
  \label{eq:lof}
\end{align}
and $c$ is the number of distinct design points, $n_i$ is the number of replicates at point $i$, $\bar{y}_{i\cdot}$ is the cell mean, and $\hat{y}_i$ is the model prediction. Under $H_0$ (no lack of fit), $F = \text{MS}_{\text{LOF}} / \text{MS}_{\text{PE}} \sim F(c - p,\; n - c)$.

We compute leave-one-out cross-validated $R^2$ using the PRESS statistic:
\begin{align}
  \text{PRESS} = \sum_{i=1}^{n} \left( \frac{e_i}{1 - h_{ii}} \right)^{\!2}, \qquad
  \text{Pred-}R^2 = 1 - \frac{\text{PRESS}}{\text{SS}_{\text{Total}}},
  \label{eq:press}
\end{align}
where $e_i$ is the $i$-th ordinary residual and $h_{ii}$ is the $i$-th diagonal element of the hat matrix. A gap between $R^2$ and Pred-$R^2$ smaller than 0.10 indicates minimal overfitting. Predicted $R^2$ and PRESS gaps vary across experiments and are reported in each experiment's model adequacy table (Tables~\ref{tab:exp1_adequacy}--\ref{tab:exp4b_adequacy}).

Gradient-boosted machines (LightGBM, 200 trees, maximum depth~4) are fit using five-fold cross-validation to establish a nonparametric upper bound on achievable $R^2$. If OLS $R^2$ is within 0.05 of LightGBM $R^2$, we conclude that the linear model with two-way interactions captures most of the signal and that higher-order or nonlinear terms contribute negligibly. In all experiments, OLS $R^2$ meets or exceeds LightGBM cross-validated $R^2$, confirming that the parametric model is well suited to the balanced factorial structure (see per-experiment diagnostics in Tables~\ref{tab:exp1_adequacy}--\ref{tab:exp4b_adequacy}).

\subsection{Variable Selection}

We fit five-fold cross-validated LASSO models with regularisation parameter $\lambda$ chosen to minimise mean squared error. Surviving variables are compared to OLS significance. In all experiments, LASSO retains auction type and number of bidders, corroborating OLS effect selection. Heredity is verified by checking that no interaction survives whose parent main effects were dropped.

The LASSO estimator \citep{Tibshirani1996} solves
\begin{align}
  \hat{\boldsymbol{\beta}}^{\text{lasso}} = \arg\min_{\boldsymbol{\beta}} \left\{ \frac{1}{2n} \|y - X\boldsymbol{\beta}\|_2^2 + \lambda \|\boldsymbol{\beta}\|_1 \right\},
  \label{eq:lasso}
\end{align}
where $\lambda \geq 0$ controls the strength of $L_1$ penalisation and is chosen by five-fold cross-validation to minimise mean squared error. The $L_1$ penalty drives small coefficients exactly to zero, performing simultaneous estimation and variable selection. Heredity requires that an interaction $\beta_{jl}$ can be nonzero only if both parent main effects $\beta_j$ and $\beta_l$ are also nonzero.

\subsection{Diagnostic Visualisations}

Four types of diagnostic plots accompany each response variable. Pareto charts rank effects by absolute $t$-statistic in descending order, with effects surpassing the critical threshold $t_{\alpha/2,\,\text{df}}$ shown in red, providing a visual hierarchy of effect importance. Main effects plots display the mean response at the low ($-1$) and high ($+1$) levels of each factor; steeper slopes correspond to larger $|\beta_j|$ values. Interaction plots show the mean response across the four combinations $(x_j, x_l) \in \{(-1,-1), (-1,+1), (+1,-1), (+1,+1)\}$ for the top six interactions by $|t|$-statistic; non-parallel lines indicate that the effect of one factor depends on the level of another. Half-normal probability plots display $|\hat{\beta}_j|$ against expected order statistics from a half-normal distribution, separating active effects (large departures from the reference line) from noise (points following the line). Residual diagnostics include quantile--quantile plots assessing normality and residuals-versus-fitted plots checking for heteroscedasticity and nonlinearity.

Under the effect sparsity principle, most factorial effects are inert and drawn from $N(0, \sigma_e^2)$; their absolute values follow a half-normal distribution. To construct the plot, the $m$ estimated effects $|\hat{\beta}_{(1)}| \leq \cdots \leq |\hat{\beta}_{(m)}|$ are sorted in ascending order and plotted against half-normal quantiles $z_i = \Phi^{-1}\!\bigl((1 + p_i)/2\bigr)$, where $p_i = (i - 0.5)/m$ \citep{Daniel1959}. Inert effects cluster along a reference line through the origin; active effects depart upward, producing a visual separation that aids identification without relying on a formal significance threshold.

\subsection{Power Analysis}

For each response variable, we compute the minimum detectable effect (MDE) at 80\% power and $\alpha = 0.05$. The standard error of each coefficient is $\hat{\sigma} / \sqrt{\sum_i x_{ij}^2}$, where $\hat{\sigma}$ is the root mean squared error and $x_{ij}$ is the $j$-th column's coded value for observation~$i$. For binary $\pm 1$ factors in a balanced design, $\sum_i x_{ij}^2 = n$, recovering the familiar $\hat{\sigma}/\sqrt{n}$ formula. For orthogonal polynomial contrasts used with three-level factors (Experiments~2 and~3), the sum of squares differs by column. The linear contrast $\{-1, 0, +1\}$ has $\sum x^2 = 2n/3$, and the quadratic contrast $\{+1, -2, +1\}$ has $\sum x^2 = 2n$. This means linear contrasts have larger standard errors (lower power) and quadratic contrasts have smaller standard errors (higher power) than binary factors at the same sample size.

All effects reported as statistically significant in the results sections survive Benjamini--Hochberg correction.
